\section{Le filtre de nouveauté en détail}\label{app:novelty}

Cette annexe donne en entier l'étape~6 du \S\ref{sec:method}, ainsi que les sources consultées pour les résultats.

\subsection{Le jeu de données d'appariements de TheoremGraph}

On commence par un passage mécanique. On résume en une phrase l'énoncé Lean du théorème réécrit et l'on cherche ce résumé dans le jeu de données d'appariements de TheoremGraph (fiche du jeu de données, Hugging Face, consultée le 23 novembre 2026) \cite{theoremgraph-matching}. Chacune de ses 100\,831 lignes associe une déclaration Lean, de Mathlib ou de l'un de 25 autres projets, à l'énoncé d'arXiv le plus proche, trouvé par similarité cosinus (au moins 0,80) entre plongements de résumés, et classé \emph{exact}, \emph{inexact} ou \emph{faux} par un modèle de langage; son juge principal tient 47\,952 paires pour des appariements. Une réponse positive établit qu'un jumeau formel existe et désigne l'article qui en donne la version informelle; le résultat va alors au journal d'étalonnage, avec une référence précise. Une réponse négative ne prouve rien: un énoncé informel sans déclaration Lean ne figure pas dans le jeu, les étiquettes sont celles d'un modèle et ne sont pas vérifiées, et le texte informel n'est fourni que pour les paires dont la source arXiv autorise la redistribution. La recherche à la main reste donc nécessaire. Nous n'avons fait ce passage que sur huit résultats (les énoncés sur les carquois, le carquois cyclique, Morita, les coins, $T_n$, $M_n(k)$, $(\mathbb Z/n,\cdot)$ et les compacts séparés; ni les automates, ni les topos combinatoires, ni les treillis de sous-espaces), par mots-clés, dans le service de recherche public du jeu, et aucune des paires renvoyées ne concernait l'un d'eux. On s'y attendait, et cela ne dit rien de la nouveauté: le côté Lean du jeu ne contient aucune déclaration de Lean Pool, ni aucune sur les topos de De Morgan, et une recherche par mots-clés classe selon les mots, non selon la similarité des résumés.

\subsection{La recherche sémantique de TheoremSearch, et l'OEIS}

Pour le côté informel, nous avons ensuite soumis les dix lignes du tableau des résultats à la recherche sémantique de TheoremSearch \cite{theoremsearch}, qui plonge une description en langue naturelle et classe quelque 8 millions d'énoncés tirés d'arXiv et d'autres sources selon la similarité de leurs résumés (API publique, 23 novembre 2026, six résultats par requête). Aucun des énoncés renvoyés ne formule l'un des dix. Trois requêtes en ont renvoyé les ingrédients, que l'article signale déjà: le théorème fondamental de la géométrie projective, pour les treillis de sous-espaces; pour le résultat sur les compacts séparés, la caractérisation des topos localiques de De Morgan (celle de Johnstone, sous la forme que lui donne Caramello) et celle des compacts séparés dont $C(X,\mathbb R)$ est complet au sens de Dedekind comme les espaces stoniens, en deux énoncés distincts; pour les coins, des énoncés sur les anneaux où une propriété de $R$ passe à $eRe$ quand $ReR = R$, mais pas l'énoncé sur les monoïdes. Les autres requêtes ont renvoyé des énoncés sur les carquois, les cycles, les monoïdes ou les diviseurs, sans rapport avec nos résultats. Avec six résultats par requête, la recherche reste peu profonde, et un énoncé formulé dans un autre vocabulaire peut lui échapper; ce passage resserre donc la recherche à la main sans la remplacer. Les appels, leurs paramètres et les réponses du 23 novembre 2026 sont dans \texttt{docs/bridges/novelty/} (\texttt{theoremsearch.py} et \texttt{theoremsearch-2026-11-23.json}): une requête \texttt{POST https://api.theoremsearch.com/search} par résultat, avec la description comme \texttt{query}, \texttt{n\_results} fixé à 6, sans filtre ni pondération par les citations. Les résultats numériques ont aussi été cherchés dans l'OEIS \cite{oeis} (23 novembre 2026): les $2^n+1$ sous-topos du $n$-cycle sont A000051; les $2^{\omega(n)}$ points essentiels des $(\mathbb Z/n,\cdot)$-ensembles sont A034444, le nombre de diviseurs unitaires de $n$, comme il se doit puisque les idempotents de $\mathbb Z/n$ correspondent à ses diviseurs unitaires; les nombres de Dedekind sont A000372; la suite $D(\omega(n))$ des sous-topos des $(\mathbb Z/n,\cdot)$-ensembles ($2, 3, 3, 3, 3, 6, 3, \dots$ pour $n=1,\dots,16$) n'est pas dans l'OEIS, bien qu'elle soit la composée de deux suites qui y sont.

\subsection{La recherche à la main}

On cherche ensuite, dans cet ordre: l'\emph{Elephant} de Johnstone; les articles de l'auteur de l'arche (ici Caramello); la littérature récente sur le tablier (ici Hemelaer, Rogers, Pirashvili pour les actions de monoïdes); la littérature classique du domaine de l'énoncé réécrit (demi-groupes, carquois, topologie); les serveurs de prépublications et le nLab. Il faut chercher l'énoncé \emph{réécrit}, dans le vocabulaire de son domaine, et pas seulement l'énoncé topos-théorique: trois de nos candidats (le critère pour les monoïdes finis, les sous-topos des $\mathbb N$-ensembles, les points essentiels comme classes d'idempotents) étaient publiés dans le langage des demi-groupes ou des actions de monoïdes, alors que nous les avions crus nouveaux. Les énoncés trouvés sont consignés au journal d'étalonnage (\S\ref{sec:calibration}).

\subsection{Sources consultées pour les résultats}

Pour chaque énoncé des \S\ref{sec:machines}--\ref{sec:shapes}, nous avons consulté l'\emph{Elephant} de Johnstone, les articles et le livre de Caramello (dont \cite{caramello-chamoun}), les articles et la thèse de Hemelaer et Rogers sur les topos d'actions de monoïdes \cite{hemelaer-rogers,rogers-1905,rogers-thesis}, la littérature sur l'équivalence de Morita et les relations de Green des monoïdes, la littérature sur les automates synchronisants, le nLab, ainsi que des recherches générales sur le web et dans les serveurs de prépublications. \q{Non trouvé} signifie exactement cela; plusieurs énoncés sont élémentaires et relèvent peut-être du folklore, et nous accueillerons volontiers des références.
